\chapter{Migration Strategies}
\label{ch:migration}

\abstract*{The post-quantum transition is not merely an algorithm swap---it requires rethinking protocols, infrastructure, and organizational processes. This chapter provides a practical framework for migrating to post-quantum cryptography, covering hybrid deployment modes, protocol updates for TLS and code signing, blockchain-specific challenges, organizational roadmaps aligned with CNSA~2.0 and European guidance, and the timeline pressures captured by Mosca's inequality.}

\abstract{The post-quantum transition is not merely an algorithm swap---it requires rethinking protocols, infrastructure, and organizational processes. This chapter provides a practical framework for migrating to post-quantum cryptography, covering hybrid deployment modes, protocol updates for TLS and code signing, blockchain-specific challenges, organizational roadmaps aligned with CNSA~2.0 and European guidance, and the timeline pressures captured by Mosca's inequality.}

% =============================================================
% SOURCE: notes.tex §5.2–5.3 (lines 6944–7876)
% REUSE: ~50%. Added CNSA 2.0, EU guidance, formal hybrid KEM
% analysis, X.509 encoding, Mosca inequality.
% =============================================================


\section{Hybrid Modes}
\label{sec:hybrid-modes}

The safest path to post-quantum security runs through \emph{hybrid cryptography}\index{hybrid cryptography}: combining a classical and a post-quantum algorithm so that the overall construction remains secure as long as \emph{either} component is secure.  This is not overcaution---it is learning from history.  The Dual~EC~DRBG episode (2006--2013) demonstrated that even standardized algorithms can conceal fatal weaknesses.  Hybrid modes protect against classical breaks of the new PQC primitive, implementation flaws in immature code, quantum timeline uncertainty, and future changes to standards.

\subsection{Composite KEMs}
\label{subsec:composite-kems}

A \emph{composite KEM}\index{KEM!composite} runs two key encapsulation mechanisms in parallel and combines their shared secrets through a key derivation function.

\begin{definition}[Composite KEM]
\label{def:composite-kem}
Let $\mathsf{KEM}_1 = (\KeyGen_1, \Encaps_1, \Decaps_1)$ and $\mathsf{KEM}_2 = (\KeyGen_2, \Encaps_2, \Decaps_2)$ be two KEMs.  The \emph{composite KEM} $\mathsf{KEM}_H$ operates as follows:
\begin{enumerate}
\item $\KeyGen_H$: Run $(pk_1, sk_1) \gets \KeyGen_1$ and $(pk_2, sk_2) \gets \KeyGen_2$.  Return $pk_H = (pk_1, pk_2)$, $sk_H = (sk_1, sk_2)$.
\item $\Encaps_H(pk_H)$: Compute $(c_1, K_1) \gets \Encaps_1(pk_1)$ and $(c_2, K_2) \gets \Encaps_2(pk_2)$.  Return $c_H = (c_1, c_2)$ and $K_H = \mathrm{KDF}(K_1 \| K_2)$.
\item $\Decaps_H(sk_H, c_H)$: Recover $K_1 \gets \Decaps_1(sk_1, c_1)$ and $K_2 \gets \Decaps_2(sk_2, c_2)$.  Return $K_H = \mathrm{KDF}(K_1 \| K_2)$.
\end{enumerate}
\end{definition}

The critical observation is that the $\mathrm{KDF}$ (e.g., HKDF-SHA256) is modelled as a random oracle: even if one shared secret is completely known to the adversary, the output remains pseudorandom as long as the other secret has sufficient entropy.

\begin{theorem}[IND-CCA2 Security of Composite KEMs]
\label{thm:hybrid-kem-security}
\index{IND-CCA2!composite KEM}
Let $\mathsf{KEM}_1$ be $(t_1, \epsilon_1)$-IND-CCA2 secure and $\mathsf{KEM}_2$ be $(t_2, \epsilon_2)$-IND-CCA2 secure.  In the random-oracle model for $\mathrm{KDF}$, the composite $\mathsf{KEM}_H$ is $(t, \epsilon)$-IND-CCA2 secure with $t = \min(t_1, t_2)$ and $\epsilon \leq \epsilon_1 + \epsilon_2$.  In particular, $\mathsf{KEM}_H$ is IND-CCA2 secure if \emph{either} component is.
\end{theorem}

\begin{svgraybox}
\textbf{Proof sketch.}  Assume an adversary $\mathcal{A}$ breaks $\mathsf{KEM}_H$ with advantage~$\epsilon$.  We construct a hybrid argument over two games: Game~0 replaces $K_1$ with a uniformly random key (reduction to $\mathsf{KEM}_1$), and Game~1 then replaces $K_2$ (reduction to $\mathsf{KEM}_2$).  In the random-oracle model, $\mathrm{KDF}(K_1 \| K_2)$ is indistinguishable from random whenever at least one input is random, so $\epsilon \leq \epsilon_1 + \epsilon_2$.  See Bindel et~al.~\cite{Bindel2019} for the full proof.
\end{svgraybox}


\subsection{Dual Signatures}
\label{subsec:dual-signatures}

For digital signatures, the hybrid approach produces two independent signatures and requires both to verify:

\begin{definition}[Composite Signature]
\label{def:composite-sig}
\index{signature!composite}
Given signature schemes $\Sigma_1 = (\KeyGen_1, \Sign_1, \Verify_1)$ and $\Sigma_2 = (\KeyGen_2, \Sign_2, \Verify_2)$, the composite scheme signs as $\sigma_H = (\Sign_1(sk_1, m),\; \Sign_2(sk_2, m))$ and accepts if and only if \emph{both} components verify.
\end{definition}

This ``AND'' composition guarantees existential unforgeability (EUF-CMA) if either component scheme is EUF-CMA secure.  The price is paid in signature and public-key size: for an RSA-2048 + ML-DSA-65 combination, the composite signature is $256 + 3{,}309 = 3{,}565$~bytes, with a composite public key of $294 + 1{,}952 = 2{,}246$~bytes.


\subsection{X.509 Hybrid Certificates}
\label{subsec:x509-hybrid}

Deploying hybrid signatures within the X.509 public-key infrastructure\index{X.509!hybrid certificates} requires encoding two public keys and two signatures into a single certificate.  Three strategies have been proposed:

\begin{enumerate}
\item \textbf{Dual-key certificate:} A single certificate carries both public keys and is signed with both CA keys.  The IETF draft \texttt{draft-ounsworth-pq-composite-sigs} defines composite OIDs for this approach.
\item \textbf{Catalyst certificate:} The classical certificate includes the PQC public key and signature in a non-critical X.509v3 extension.  Legacy clients ignore the extension; PQC-aware clients verify both.
\item \textbf{Paired certificates:} Two separate certificates---one classical, one PQC---are bound by a shared serial number.  The server sends both; the client verifies whichever it supports.
\end{enumerate}

The catalyst approach offers the best backward compatibility: existing TLS stacks parse the certificate without modification, while updated clients gain quantum resistance.  The cost is roughly $5$--$6$~KB of additional data per certificate (dominated by the ML-DSA public key and CA signature).


\subsection{Performance Impact}
\label{subsec:hybrid-performance}

\begin{table}[t]
\caption{Hybrid KEM overhead: X25519 + ML-KEM-768}
\label{tab:hybrid-overhead}
\begin{tabular}{p{3.5cm}p{2cm}p{2cm}p{2cm}p{2cm}}
\hline\noalign{\smallskip}
\textbf{Operation} & \textbf{X25519} & \textbf{ML-KEM-768} & \textbf{Hybrid} & \textbf{Factor} \\
\noalign{\smallskip}\svhline\noalign{\smallskip}
Key generation & 45\,\textmu s & 35\,\textmu s & 80\,\textmu s & 1.8$\times$ \\
Encapsulation   & 45\,\textmu s & 44\,\textmu s & 89\,\textmu s & 2.0$\times$ \\
Decapsulation   & 45\,\textmu s & 40\,\textmu s & 85\,\textmu s & 1.9$\times$ \\
Ciphertext size  & 32\,B & 1{,}088\,B & 1{,}120\,B & 35$\times$ \\
\noalign{\smallskip}\hline
\end{tabular}
\end{table}

The computational overhead of hybrid KEMs is modest ($< 2\times$), because ML-KEM operations are fast.  The dominant cost is \emph{bandwidth}: the combined ciphertext and public-key material adds roughly 2\,KB per handshake, which matters on constrained links but is negligible on broadband.


%% ================================================================
\section{TLS Migration}
\label{sec:tls-migration}

TLS~1.3 protects billions of HTTPS connections daily and is the highest-priority migration target.  Its clean key-exchange design---a single round-trip with ephemeral Diffie--Hellman---makes post-quantum integration relatively straightforward compared with older protocols.

\subsection{TLS 1.3 with Post-Quantum Key Exchange}
\label{subsec:tls13-pq}

TLS~1.3 negotiates key exchange via \emph{named groups}\index{TLS~1.3!named groups} advertised in the \texttt{supported\_groups} extension.  Adding a post-quantum or hybrid KEM requires:
\begin{enumerate}
\item Assigning a new named-group code point (e.g., \texttt{0x6399} for X25519Kyber768\index{X25519Kyber768}).
\item Extending buffer sizes: the \texttt{ClientHello} key-share grows from 32~bytes (X25519) to $\approx$\,1{,}216~bytes; the \texttt{ServerHello} encapsulation adds $\approx$\,1{,}088~bytes.
\item Combining the shared secrets via the TLS~1.3 key schedule: $\mathit{SS}_{\mathrm{hybrid}} = \mathrm{HKDF\text{-}Extract}(K_{\mathrm{X25519}} \| K_{\mathrm{ML\text{-}KEM}})$.
\end{enumerate}

Crucially, the handshake \emph{message count} does not change---the post-quantum material is embedded in the existing \texttt{key\_share} extension---so the protocol retains its 1-RTT structure.

\subsection{Chrome and Cloudflare Experiments}
\label{subsec:chrome-cloudflare}

\begin{example}[CECPQ2 and Beyond]
\label{ex:cecpq2}
\index{CECPQ2}
Google and Cloudflare have conducted the largest post-quantum TLS experiments to date:
\begin{itemize}
\item \textbf{2019:} CECPQ2 deployed hybrid ECDH + NTRU-HRSS in Chrome, covering $\sim$30\% of TLS connections to Google servers.  Median latency increase: 1\,ms.  User complaints: zero.
\item \textbf{2022:} Cloudflare switched to hybrid X25519 + Kyber-768, making PQC available across 20+ million websites.
\item \textbf{2023--2024:} Chrome enabled X25519Kyber768 by default.  By late 2024, over 10\% of all global HTTPS traffic used a post-quantum key exchange.
\end{itemize}
These experiments demonstrated that Internet-scale hybrid deployment is feasible with negligible user-facing impact.
\end{example}

\subsection{X25519Kyber768}
\label{subsec:x25519kyber768}

X25519Kyber768 (also known as \texttt{X25519MLKEM768} after the FIPS~203~\cite{FIPS203} renaming) is the first hybrid KEM deployed at scale.  It concatenates the 32-byte X25519 shared secret with the 32-byte ML-KEM-768 shared secret and feeds the 64-byte result into the TLS~1.3 HKDF key schedule.  The combined ciphertext is $32 + 1{,}088 = 1{,}120$~bytes---well within a single TCP segment, avoiding handshake fragmentation in most network conditions.


\subsection{Certificate Chain Sizes and Latency}
\label{subsec:cert-chain-sizes}

While key exchange adds $\sim$2\,KB, the larger concern is \emph{certificate chains}\index{certificate chain!PQC overhead}.  A three-certificate chain with ML-DSA-65 signatures grows from $\sim$3\,KB (RSA-2048) to $\sim$17\,KB, a 5.5$\times$ increase.

\begin{table}[t]
\caption{TLS handshake size and latency overhead}
\label{tab:tls-overhead}
\begin{tabular}{p{4cm}p{2.5cm}p{2.5cm}p{2.5cm}}
\hline\noalign{\smallskip}
\textbf{Component} & \textbf{Classical} & \textbf{PQC/Hybrid} & \textbf{Increase} \\
\noalign{\smallskip}\svhline\noalign{\smallskip}
Key exchange & 96\,B & 2{,}336\,B & 24$\times$ \\
Certificate chain (3 certs) & 3{,}150\,B & 17{,}235\,B & 5.5$\times$ \\
Total handshake & $\sim$3.5\,KB & $\sim$20\,KB & 5.7$\times$ \\
Median latency & baseline & +1.3\,ms & --- \\
95th-pct latency & baseline & +21\,ms & --- \\
\noalign{\smallskip}\hline
\end{tabular}
\end{table}

On desktop broadband, the extra latency is invisible.  On mobile networks with high round-trip times, the 95th-percentile overhead can reach 20--90\,ms, primarily due to TCP-level fragmentation and retransmission when the handshake exceeds the initial congestion window.  Session resumption via TLS~1.3 PSK mode avoids this cost entirely and becomes \emph{more important} in a post-quantum world.

\begin{svgraybox}
\textbf{The middlebox problem.}  Corporate firewalls and deep-packet-inspection appliances often impose hard limits on TLS record sizes (commonly 8\,KB or 16\,KB).  PQC certificate chains that exceed these limits are silently dropped, causing mysterious connection failures.  Early deployers report 2--6~weeks of debugging before identifying the root cause.  The fix is to fragment certificates across multiple TLS records or to use certificate compression (RFC~8879).
\end{svgraybox}


%% ================================================================
\section{Code Signing and Software Supply Chain}
\label{sec:code-signing}

Code signing\index{code signing} presents challenges distinct from TLS: signatures are verified \emph{years} after creation, backward compatibility with legacy verifiers is critical, and signature size directly affects software distribution costs.

\subsection{Firmware Updates}
\label{subsec:firmware}

Embedded devices---IoT sensors, medical implants, industrial controllers---verify firmware signatures using public keys burned into hardware at manufacture.  These keys may need to remain valid for 10--20~years, making them high-priority targets for ``harvest now, decrypt later'' attacks on confidentiality and, more critically, for future signature-forgery attacks on integrity.  A quantum adversary who can forge a firmware signature can push malicious updates to every device in the fleet.

Stateful hash-based signature schemes (LMS\index{LMS}, XMSS\index{XMSS}), standardized in NIST~SP~800-208, are attractive for firmware signing because their security rests purely on hash-function properties and they produce compact signatures ($\sim$2.5\,KB for LMS with $H=20$).  The main operational burden is state management: each private key must track which one-time keys have been used, and reuse is catastrophic.  For CA root keys that sign infrequently, this burden is manageable.

\subsection{Package Managers and Repository Signing}
\label{subsec:package-managers}

Software supply chains (npm, PyPI, apt, Homebrew) sign repository metadata and individual packages.  A single compromised signing key can inject malicious code into millions of systems.  Migration considerations include:
\begin{itemize}
\item \textbf{Dual signatures:} Repositories can distribute both classical and PQC signatures in parallel.  Clients verify whichever their toolchain supports.
\item \textbf{Transparency logs:} Systems like Sigstore already decouple signatures from distribution, easing algorithm transitions.
\item \textbf{Threshold signing:} Multi-party signing ceremonies (common for root keys) must be updated to use PQC-compatible threshold protocols.
\end{itemize}

\subsection{The Signature Size Challenge}
\label{subsec:sig-size-challenge}

\begin{table}[t]
\caption{Signature sizes for code-signing candidates}
\label{tab:codesign-sizes}
\begin{tabular}{p{3.5cm}p{2.5cm}p{2.5cm}p{3cm}}
\hline\noalign{\smallskip}
\textbf{Scheme} & \textbf{Signature} & \textbf{Public key} & \textbf{Security basis} \\
\noalign{\smallskip}\svhline\noalign{\smallskip}
RSA-2048 & 256\,B & 294\,B & Factoring \\
ECDSA-P256 & 64\,B & 64\,B & ECDLP \\
ML-DSA-65 & 3{,}309\,B & 1{,}952\,B & Module-LWE \\
SLH-DSA-128s & 7{,}856\,B & 32\,B & Hash functions \\
LMS ($H{=}20$) & 2{,}508\,B & 56\,B & Hash functions \\
\noalign{\smallskip}\hline
\end{tabular}
\end{table}

For a typical software update containing 1{,}000--10{,}000 signed files, switching from RSA-2048 to ML-DSA-65 increases total signature overhead from $\sim$250\,KB to $\sim$3.3\,MB per update.  SLH-DSA offers the advantage of minimal public-key size and conservative, hash-only security assumptions, but its signatures are 2--3$\times$ larger than ML-DSA.  The practical choice depends on whether the deployment is bandwidth-constrained (favour ML-DSA) or demands maximal cryptographic conservatism (favour SLH-DSA or LMS for root keys).

The signature-size tension is annoying for software distribution; for blockchain systems, where every byte is replicated across thousands of nodes and stored forever, it becomes an existential design constraint.

%% ================================================================
\section{Blockchain and Cryptocurrencies}
\label{sec:blockchain}

Blockchain systems face a ``perfect storm'' of post-quantum migration challenges: immutable transaction histories, size-sensitive economics, decentralized governance, and catastrophic failure modes.

\subsection{ECDSA Vulnerability}
\label{subsec:ecdsa-vulnerability}

Bitcoin, Ethereum, and most cryptocurrencies derive address security from ECDSA over the \texttt{secp256k1} curve.  Shor's algorithm\index{Shor's algorithm!ECDSA} recovers private keys from public keys in polynomial time.  In Bitcoin, the public key is exposed whenever coins are \emph{spent} from a pay-to-public-key-hash (P2PKH) address.  An adversary with a cryptographically relevant quantum computer (CRQC) could:
\begin{enumerate}
\item Monitor the mempool for pending transactions that reveal a public key.
\item Extract the private key before the transaction is confirmed.
\item Submit a competing transaction redirecting the funds.
\end{enumerate}

More dangerously, coins held in addresses where the public key was previously exposed (estimated at over 4~million BTC, roughly 20\% of all supply) are vulnerable even without monitoring the mempool.

\subsection{Migration Challenges}
\label{subsec:blockchain-migration}

The ECDSA vulnerability is clear enough; what makes blockchain migration uniquely difficult is that every proposed remedy collides with a core design property.

Consider immutability first.  The blockchain's defining feature---that no past record can be altered---becomes a liability when past records contain quantum-vulnerable signatures.  Historical ECDSA signatures cannot be replaced, only supplemented with new PQC protections going forward.  This stands in sharp contrast to a TLS deployment, where rotating to new certificates immediately protects all future sessions and the old certificates simply expire.

Transaction size compounds the problem.  Replacing ECDSA ($\sim$72\,B per signature) with ML-DSA-44 ($\sim$2{,}420\,B) inflates a simple Bitcoin transaction from $\sim$250\,B to $\sim$3{,}500\,B---a 14$\times$ increase.  Block capacity drops from $\sim$4{,}000 to $\sim$285~transactions, fees rise proportionally, and full-node storage grows from $\sim$50\,GB/year to $\sim$700\,GB/year.  In a system where every full node stores the complete history, this is not merely a bandwidth cost; it reshapes the economics of participation.

Governance adds a further constraint.  Any signature-scheme change requires a hard fork or, at minimum, a soft fork with broad community agreement.  In decentralized networks, achieving such agreement can take years---the Bitcoin block-size debate (2015--2017) is a cautionary precedent.

Finally, address migration demands active participation from every user.  Funds must be moved to quantum-safe addresses through on-chain transactions, which means every wallet holder must act.  Funds in dormant wallets---including those belonging to lost keys, estimated at 3--4~million BTC---remain permanently vulnerable, creating a systemic risk that no protocol upgrade can eliminate unilaterally.

\subsection{Frozen Funds Risk}
\label{subsec:frozen-funds}

\index{frozen funds risk}
Coins whose owners cannot or will not migrate before a CRQC arrives face outright theft.  The community has proposed several mitigation strategies, none of them painless.

The most conservative approach is \emph{commit-reveal}: users commit a hash of a PQC public key to the blockchain now, then reveal the actual key only when spending.  Because the quantum-vulnerable classical public key never enters the mempool, a CRQC operator cannot intercept and forge a competing transaction.  The limitation is clear---commit-reveal protects future transactions but does nothing for the millions of addresses whose public keys were exposed in past spends.

A second line of defence is the \emph{quantum-safe sidechain}: a parallel chain secured by PQC signatures that accepts bridged assets from the main chain.  Users lock funds on the classical chain and receive equivalent tokens on the sidechain, migrating at their own pace without requiring a hard fork.  The approach is technically sound but adds operational complexity and splits network effects between two chains.

The most drastic proposal is \emph{deadline-based freezing}.  After a community-agreed date, the network would reject transactions from addresses that have not migrated to PQC keys.  This protects the system from quantum theft but permanently locks affected funds---including coins belonging to deceased holders or lost wallets.  The measure is politically explosive; it amounts to a collective decision to sacrifice some holders' assets for the security of the whole, a trade-off with no precedent in Bitcoin's history.

\begin{example}[Ethereum's Account Abstraction Path]
\label{ex:ethereum-pqc}
Ethereum's account-abstraction roadmap (EIP-4337 and successors) decouples signature verification from the protocol layer, allowing smart-contract wallets to adopt PQC signature schemes without a hard fork.  The planned sequence is: enable smart-contract wallets (2024--2025), support hybrid signatures (2025--2027), incentivize migration via gas discounts (2027--2029), and disable raw ECDSA for new transactions (2030+).  This is the most credible blockchain migration plan to date, but it still requires years of coordinated effort.
\end{example}

Blockchain migration illustrates, in extreme form, the tension between technical necessity and organizational inertia.  The next section generalizes this tension into a structured roadmap applicable to any organization, blockchain or otherwise.


%% ================================================================
\section{Organizational Roadmap}
\label{sec:roadmap}

The preceding sections examined specific migration targets---TLS, code signing, blockchain---each with its own technical constraints.  But even perfect technical solutions fail without organizational execution.  Organizations must run a structured migration program that touches procurement, engineering, compliance, and operations in coordinated sequence.  The six-phase roadmap below distills guidance from NIST~SP~1800-38, the NSA's CNSA~2.0 advisory, and European agency recommendations into a practical framework.  It is ordered by dependency: each phase produces the inputs the next phase requires.

\subsection{Phase 1: Cryptographic Inventory}
\label{subsec:phase-inventory}

\index{cryptographic inventory}
No migration can succeed without knowing what must be moved.  The first phase is a comprehensive discovery of every use of public-key cryptography across the organization: TLS certificates, code-signing keys, VPN configurations, SSH keys, API tokens, hardware security modules, embedded-device firmware, and third-party dependencies.  The analogy to physical logistics is apt---before relocating a warehouse, one inventories every pallet; skipping this step guarantees that critical items are left behind.

Automated scanning tools (Venafi, Keyfactor, or open-source \texttt{cbom} generators) provide a starting point, but they rarely find everything.  Manual review is essential, particularly for shadow IT and legacy systems maintained by teams outside central security.  A typical medium enterprise discovers that 40--50\% of its cryptographic assets were unknown to the security team before the inventory---a sobering figure that explains why so many migrations stall in later phases.

With a complete inventory in hand, the organization can move to risk assessment: not all assets face the same quantum exposure, and treating them uniformly wastes resources.

\subsection{Phase 2: Risk Assessment}
\label{subsec:phase-risk}

The inventory tells us \emph{what} the organization has; risk assessment determines \emph{what is at stake}.  Each asset is classified along two dimensions: \emph{data sensitivity lifetime}---how long must the data remain confidential or its signatures remain trusted?---and \emph{migration difficulty}---how hard is it to change the algorithm?  Assets that score high on both dimensions (long-lived secrets protected by deeply embedded cryptography) demand immediate attention; ephemeral session keys on modern, patchable servers can wait.

\begin{table}[t]
\caption{Risk classification matrix}
\label{tab:risk-matrix}
\begin{tabular}{p{3.5cm}p{2.5cm}p{2.5cm}p{3cm}}
\hline\noalign{\smallskip}
\textbf{Data category} & \textbf{Lifetime} & \textbf{Risk level} & \textbf{Action} \\
\noalign{\smallskip}\svhline\noalign{\smallskip}
Government / military & 25--50 years & Critical & Immediate \\
Healthcare records & 50--100 years & Critical & Immediate \\
Financial records & 7--10 years & High & By 2027 \\
Trade secrets / IP & 10--20 years & High & By 2027 \\
Personal data (GDPR) & Lifetime & Medium & By 2029 \\
Ephemeral session keys & Hours & Low & By 2031 \\
\noalign{\smallskip}\hline
\end{tabular}
\end{table}

\subsection{Phase 3: Prioritize}
\label{subsec:phase-prioritize}

Risk assessment produces a ranked list; prioritization turns that list into a work plan.  The principle is straightforward---start with the highest-risk, longest-lived data---but the practical ordering requires judgment about dependencies and organizational capacity.

VPN tunnels and TLS sessions protecting classified or regulated data typically top the list, because they are actively exposed to harvest-now-decrypt-later attacks and Mosca's inequality (Sect.~\ref{subsec:mosca}) is already violated for multi-decade secrets.  Root CA and code-signing keys follow closely: they are long-lived, their compromise has cascading impact, and their migration is largely independent of application-layer changes.  Data-at-rest encryption keys for archives with multi-decade retention come next, followed by customer-facing web services, which benefit from well-understood migration paths (Sect.~\ref{sec:tls-migration}) and high public visibility.

The prioritized list feeds directly into the testing phase, where each migration path is validated before it touches production traffic.

\subsection{Phase 4: Test}
\label{subsec:phase-test}

Before any PQC algorithm touches production traffic, it must survive contact with reality.  Hybrid modes are deployed first in staging and pre-production environments, where failures are cheap and reversible.

Testing spans several dimensions.  Functional correctness is the baseline: do the hybrid handshakes complete, and do the resulting session keys match on both sides?  Latency and throughput under load reveal whether the larger key material degrades performance beyond acceptable thresholds---the middlebox problem described in Sect.~\ref{subsec:cert-chain-sizes} is routinely discovered at this stage.  Interoperability testing with legacy clients and middleboxes is critical, because a migration that breaks 2\% of connections will be rolled back before it protects anyone.  Failure-mode behavior must be explicitly validated: does the system fall back gracefully to classical-only when the PQC component is unavailable, or does it fail hard?  Finally, key and certificate lifecycle operations---renewal, revocation, rotation---must work end-to-end, because a migration that cannot maintain itself is a migration that will decay.

Testing failures are not setbacks; they are the point.  Every incompatibility caught here is one that does not become a 2\,AM incident in production.

\subsection{Phase 5: Deploy}
\label{subsec:phase-deploy}

Deployment follows the same graduated-rollout discipline used for any high-risk infrastructure change---the difference is that the stakes include long-term cryptographic exposure, not just transient outages.

The rollout begins with a canary deployment to roughly 1\% of traffic, chosen to include a representative mix of client populations.  At this scale, gross incompatibilities surface---broken middleboxes, misconfigured load balancers, clients that choke on larger handshakes---without affecting the majority of users.  Once the canary is stable, traffic shifts to 10\% (early adopters), where performance under realistic load is validated against the baselines established in testing.  The move to 50\% exercises the system at scale: tail latencies, connection-retry rates, and certificate-chain failures that appear only under heavy load become visible.  Full deployment at 100\% completes the migration, but the classical algorithm is not removed---it remains available as a fallback until the organization is confident that every client ecosystem has adapted.

Throughout each stage, monitoring dashboards track handshake failure rates, latency distributions, and certificate-validation errors.  Any regression beyond predefined thresholds triggers an automatic rollback to the previous traffic split.

\subsection{Phase 6: Verify and Maintain}
\label{subsec:phase-verify}

A migration that declares victory at 100\% deployment and moves on has failed.  Cryptographic infrastructure is a living system: new services are provisioned, dependencies are updated, developers under deadline pressure reach for the defaults their frameworks provide---and those defaults may still be classical.

Continuous compliance therefore requires several ongoing disciplines.  Monitoring detects \emph{algorithm drift}, where newly deployed services silently revert to classical-only configurations.  Automated alerting fires when certificates or keys use deprecated algorithms, catching regressions before they reach production.  Periodic reassessment accounts for advances in cryptanalysis: a parameter set considered safe in 2026 may need revision by 2030 as lattice-reduction techniques improve.

The deepest lesson of this phase is the importance of \emph{crypto-agility}\index{crypto-agility}---the architectural ability to swap algorithms without redesigning protocols or redeploying hardware.  Crypto-agility is not a feature to add later; it is a design principle that must be embedded from Phase~1 onward.  Organizations that achieve it will navigate not only the post-quantum transition but every subsequent cryptographic migration with far less friction.

The CNSA~2.0 timeline and European guidance that follow provide the external deadlines against which these six phases must be scheduled.

\subsection{CNSA 2.0 Timeline}
\label{subsec:cnsa20}

\index{CNSA 2.0}
The six phases above describe \emph{how} to migrate; the CNSA~2.0 timeline specifies \emph{when}.  The NSA's Commercial National Security Algorithm Suite~2.0 sets mandatory deadlines for U.S.\ national-security systems, and these deadlines increasingly serve as de facto benchmarks for the private sector as well:

\begin{table}[t]
\caption{CNSA 2.0 transition deadlines (NSA, 2022)}
\label{tab:cnsa20}
\begin{tabular}{p{4.5cm}p{3.5cm}p{3.5cm}}
\hline\noalign{\smallskip}
\textbf{Use case} & \textbf{Algorithm} & \textbf{Deadline} \\
\noalign{\smallskip}\svhline\noalign{\smallskip}
Software / firmware signing & CNSA~2.0 exclusive & 2025 \\
Web servers / cloud services & CNSA~2.0 preferred & 2025 \\
Traditional networking (VPN, routers) & CNSA~2.0 preferred & 2026 \\
Operating systems & CNSA~2.0 support & 2027 \\
Custom / embedded applications & CNSA~2.0 exclusive & 2030 \\
Legacy infrastructure & CNSA~2.0 exclusive & 2033 \\
\noalign{\smallskip}\hline
\end{tabular}
\end{table}

\subsection{European Guidance}
\label{subsec:eu-guidance}

CNSA~2.0 is a U.S.-centric framework; European agencies set their own timelines, informed by different threat models and regulatory structures, but converging on similar urgency.

ANSSI (France) recommends hybrid modes as mandatory during the transition period and requires PQC for new systems handling classified data from 2025 onward.  Notably, ANSSI explicitly endorses FrodoKEM alongside ML-KEM as a conservative alternative for high-security applications---reflecting a preference for schemes with stronger theoretical security reductions, even at a performance cost.  BSI (Germany) recommends beginning migration planning immediately, with hybrid deployments by 2025 and full PQC by 2030; BSI requires a ``defence in depth'' approach combining lattice-based and hash-based schemes, ensuring that no single cryptanalytic breakthrough compromises the entire infrastructure.  At the EU level, ENISA published a post-quantum roadmap in 2022 urging member states to conduct cryptographic inventories by 2025 and complete critical-infrastructure migration by 2030.

The convergence across these agencies---American, French, German, European---is striking.  Despite differing institutional cultures and threat priorities, all place the start of serious migration activity no later than 2025 and demand completion for critical systems by 2030.  Whether these deadlines are achievable depends on the timeline pressures we examine next.


%% ================================================================
\section{Timeline and Urgency}
\label{sec:timeline}

The roadmap and regulatory deadlines of the previous section raise an obvious question: how late is too late?  The answer depends not on policy but on arithmetic---a simple inequality that formalizes the harvest-now-decrypt-later threat into a planning tool.

\subsection{Mosca's Inequality}
\label{subsec:mosca}

\index{Mosca's inequality}
Michele Mosca formalized the urgency of post-quantum migration with a simple but powerful inequality:

\begin{definition}[Mosca's Inequality]
\label{def:mosca}
Let:
\begin{itemize}
\item $x$ = the number of years the data must remain secure,
\item $y$ = the number of years needed to migrate the system to PQC,
\item $z$ = the number of years until a CRQC exists.
\end{itemize}
If $x + y > z$, then the data is \emph{already at risk today}, because an adversary can harvest ciphertext now and decrypt it once a CRQC is available before the confidentiality requirement expires.
\end{definition}

\begin{example}[Applying Mosca's Inequality]
\label{ex:mosca-hospital}
A hospital system stores patient records that must remain confidential for 50~years ($x = 50$).  Their IT department estimates that a full cryptographic migration will take 5~years ($y = 5$).  Assuming a CRQC arrives by 2035 ($z = 2035 - 2026 = 9$), we have $x + y = 55 \gg 9 = z$.  The inequality has been violated for decades: any patient record encrypted before 2026 with a classical algorithm is exposed to a harvest-now-decrypt-later attack.
\end{example}

The uncomfortable implication: for any data with a multi-decade confidentiality requirement, migration should have started years ago.

\subsection{Sector-Specific Deadlines}
\label{subsec:sector-deadlines}

Mosca's inequality is most useful when applied to concrete sectors, because different industries face radically different values of $x$, $y$, and~$z$:

\begin{table}[t]
\caption{Sector-specific Mosca analysis (assuming CRQC by 2035)}
\label{tab:sector-mosca}
\begin{tabular}{p{3cm}p{1.5cm}p{1.5cm}p{1.5cm}p{3.5cm}}
\hline\noalign{\smallskip}
\textbf{Sector} & $x$ (yr) & $y$ (yr) & $x+y$ & \textbf{Must start by} \\
\noalign{\smallskip}\svhline\noalign{\smallskip}
Intelligence & 50 & 8 & 58 & Already overdue \\
Healthcare & 50 & 5 & 55 & Already overdue \\
Finance & 10 & 4 & 14 & 2026 \\
Telecoms & 5 & 3 & 8 & 2027 \\
Retail / web & 2 & 2 & 4 & 2031 \\
\noalign{\smallskip}\hline
\end{tabular}
\end{table}

\subsection{Cost of Delay}
\label{subsec:cost-delay}

Delaying migration incurs compounding costs:
\begin{enumerate}
\item \textbf{Harvest-now exposure grows daily.}  Every additional day of classical-only encryption adds one more day of ciphertext to the adversary's stockpile.
\item \textbf{Technical debt accumulates.}  New systems deployed today without PQC support become tomorrow's ``legacy infrastructure''---the most expensive category to migrate.
\item \textbf{Talent scarcity.}  PQC expertise is limited.  Organizations that wait will compete for the same engineers during a compressed timeline, driving up costs.
\item \textbf{Regulatory risk.}  Once mandates like CNSA~2.0 take effect, non-compliant systems face operational restrictions or loss of certification.
\end{enumerate}

\begin{svgraybox}
\textbf{Lessons from previous cryptographic migrations.}
The AES transition (2001--2010) succeeded because it offered a drop-in replacement with better performance.  The IPv6 transition (1998--present, still incomplete) stalled because it required infrastructure changes with no immediate benefit.  PQC migration resembles IPv6 more than AES: the threat is not yet tangible, the changes are invasive, and the performance trade-offs are unfavorable.  Without deliberate urgency, a 20-year transition is a realistic worst case.
\end{svgraybox}


%% ================================================================
\section*{Problems}
\addcontentsline{toc}{section}{Problems}

\begin{prob}
\label{prob:ch13-mosca}
Apply Mosca's inequality to the following scenario: a hospital system stores patient records that must remain confidential for 50~years.  Their IT department estimates that a full cryptographic migration will take 5~years.  Assuming a cryptographically relevant quantum computer arrives in 2035:
\begin{enumerate}
\item[(a)] By what year must the migration begin?
\item[(b)] Repeat the analysis for a retail company whose transaction data must remain confidential for only 2~years, with a migration time of 18~months.
\item[(c)] In which case is the urgency greater, and why?
\end{enumerate}
\end{prob}

\begin{prob}
\label{prob:ch13-hybrid-kem}
Design a hybrid KEM that combines X25519 and ML-KEM-768.  Specify how the shared secrets are combined (concatenation followed by a KDF), and argue---using the structure of Theorem~\ref{thm:hybrid-kem-security}---that your construction is IND-CCA2 secure if either component KEM is secure.
\end{prob}

\begin{prob}
\label{prob:ch13-cert-chain}
A TLS server presents a certificate chain of depth~3 (end-entity, intermediate CA, root CA).  Compute the total chain size in bytes for each of the following configurations:
\begin{enumerate}
\item[(a)] All certificates use RSA-2048 (public key: 294\,B, signature: 256\,B, overhead: 200\,B per cert).
\item[(b)] All certificates use ML-DSA-65 (public key: 1{,}952\,B, signature: 3{,}309\,B, overhead: 200\,B per cert).
\item[(c)] A hybrid catalyst certificate (RSA public key + ML-DSA public key and signature in an extension).
\end{enumerate}
Discuss which configuration exceeds a typical TCP initial congestion window of 14{,}600~bytes and what mitigations are available.
\end{prob}

\begin{prob}
\label{prob:ch13-blockchain}
Bitcoin transactions currently average 250~bytes with an ECDSA signature of 72~bytes.  Suppose ECDSA is replaced by ML-DSA-44 (signature: 2{,}420\,B; public key: 1{,}312\,B).
\begin{enumerate}
\item[(a)] Compute the new average transaction size.
\item[(b)] If the block size limit remains 4\,MB, estimate the new maximum number of transactions per block.
\item[(c)] Discuss two alternative approaches that could mitigate the throughput reduction.
\end{enumerate}
\end{prob}

\begin{prob}
\label{prob:ch13-roadmap}
An enterprise operates 3{,}000 TLS endpoints, 200 code-signing certificates, and 50 VPN gateways.  Using the six-phase roadmap of Sect.~\ref{sec:roadmap}:
\begin{enumerate}
\item[(a)] Identify which assets should be migrated first and justify your prioritization using Mosca's inequality.
\item[(b)] Propose a 4-year timeline with milestones for each phase.
\item[(c)] Explain why crypto-agility (Phase~6) is essential even after migration is ``complete.''
\end{enumerate}
\end{prob}
